\section{Conclusion}
\label{sec:conclusion}

This report has documented the group engineering process behind the AENGM0074 autonomous SAR drone project. A team of five MSc students collaborated over 12 weeks to design, build, and test a quadcopter system capable of autonomous search, AI-based casualty detection, operator-in-the-loop verification, and offset landing.

Our key achievements as a team were:

\begin{itemize}
  \item \textbf{Effective parallel development}: five workstreams running simultaneously with clear interfaces and integration checkpoints, enabled by a modular software architecture.
  \item \textbf{Simulation-first validation}: the complete system was validated in simulation before hardware integration, dramatically reducing integration risk and debugging time.
  \item \textbf{Progressive test methodology}: a five-level test progression ensured that no untested capability was introduced to a live aircraft, with over 40 purpose-built test scripts providing comprehensive coverage.
  \item \textbf{Structured decision-making}: multi-criteria trade studies for key design choices (detection model, streaming protocol, companion computer) documented in a persistent design decisions log.
  \item \textbf{Proactive risk management}: a risk register reviewed at each sprint, with cross-cutting mitigations (simulation-first, multiple model variants, dual communication paths) that addressed multiple risks simultaneously.
  \item \textbf{Comprehensive documentation}: architecture documents, setup guides, flight checklists, and blueprints version-controlled alongside the code, ensuring knowledge persistence across the team.
\end{itemize}

The project also highlighted areas for improvement, particularly around earlier hardware integration, more structured sprint ceremonies, deeper cross-training, and booking sufficient outdoor testing time. These lessons are offered to future teams undertaking similar projects.

Despite the constraints of a compressed academic timeline and weather-limited field access, the team delivered a fully integrated system that was validated through bench testing, video analysis of real flight footage, and calibrated hardware measurements. The autonomous SAR drone project demonstrated that effective interdisciplinary collaboration, structured systems engineering, and disciplined progressive testing can deliver a complex autonomous system within the tight constraints of a university group project.
